\documentclass[10pt]{article}

% -----------------------------------------------------------------------
% TMLR paper (DRAFT): A Benchmark Recipe Is Not an Experiment
% Target: Transactions on Machine Learning Research (TMLR)
%
% This is a scaffolded first draft generated from the master collaborative
% reference (docs/Master_Collaborative_Reference_2026-07-15c/). Content that
% depends on experiments not yet run is wrapped in a boxed "[SCAFFOLD --- to be
% completed]" block or marked inline with "[PENDING: ...]"; every such marker is
% a concrete TODO for the empirical program. Grep for SCAFFOLD and pending: to
% enumerate them.
%
% COMPILE: drops in the official TMLR style if tmlr.sty is present; otherwise
% falls back to an article layout so the draft compiles as-is.
% -----------------------------------------------------------------------

\IfFileExists{tmlr.sty}{%
  \usepackage{tmlr}%
}{%
  % --- fallback so the draft compiles without the TMLR class ---
  \usepackage[margin=1in]{geometry}%
  \usepackage{fancyhdr}%
  \usepackage{natbib}%          % tmlr.sty provides natbib itself; load only in fallback
  \pagestyle{fancy}\fancyhf{}\fancyhead[C]{\small\itshape Draft --- not for distribution}\fancyfoot[C]{\thepage}%
}

\usepackage[T1]{fontenc}
\usepackage[utf8]{inputenc}
\usepackage{microtype}
\usepackage{amsmath,amssymb,amsfonts}
\usepackage{booktabs}
\usepackage{tabularx}
\usepackage{array}
\usepackage{float}
\usepackage{xcolor}
\usepackage{enumitem}
\usepackage{graphicx}
\usepackage[colorlinks=true,linkcolor=black,citecolor=black,urlcolor=blue]{hyperref}
\usepackage{url}
\usepackage{xurl}

% ---- conventions (no custom macros: everything below is written out at each
% use site, so any single file compiles against any copy of this preamble) ----
%
% Inline code and file names: \texttt{...}.
%
% Draft markers: a scaffold block is an \fcolorbox{orange}{yellow!12} around a
% 0.97\linewidth minipage headed "[SCAFFOLD --- to be completed]"; a pending
% marker is {\color{orange}\textsf{\small[\textsc{pending}: ...]}}.
%
% tabularx tables weight their columns inline, with
% >{\raggedright\arraybackslash\hsize=<w>\hsize}X. Do not go back to mixing
% fixed p{<f>\linewidth} columns with X: tabularx sizes the table to \linewidth
% either way, but a justified narrow p{} column lets a word it cannot hyphenate
% protrude into the margin. The weights of one table must sum to its number of
% columns.

\title{A Benchmark Recipe Is Not an Experiment:\\
Attribution and Identifiability in the Reproduction of Open-Weight LLM Evaluations}

\author{%
  Edward Wang\thanks{Kitware, Inc. Correspondence: \texttt{edward.wang@kitware.com}.}
  \quad\textnormal{{\color{orange}\textsf{\small[\textsc{pending}: co-authors / advisor]}}}\\
  Kitware, Inc.
}

\date{Draft compiled \today}

\begin{document}
\maketitle

\begin{abstract}
Public large-language-model benchmarks such as HELM publish, for each
(scenario, model) pair, the prompts sent, the completions returned, and the
computed metrics. It is tempting to treat these artifacts as a complete
experimental record and to ask, simply, whether a locally re-run evaluation
reproduces the published number. We argue --- and demonstrate on the HELM
corpus --- that this framing is ill-posed: a benchmark \emph{recipe} (scenario,
prompt construction, decoding, scorer) does not identify the \emph{execution
instrument} that produced a score. The surviving record routinely omits the
model and tokenizer revision, the load precision, the chat-template version, the
serving engine and its kernels, hardware-sensitive numerical settings, the
harness version, and the artifact-migration history. Reproduction is therefore
better cast as a layered \emph{identification} problem. We present
\textsc{EvalAudit}, a system for exact recipe replay, execution-substrate
reconstruction, version-pinned (``era'') scoring instruments, and layered
differential comparison, and we use it to (i) attribute concrete disagreements on
modern open-weight models (OLMo) to specific unrecorded parameters --- a tokenizer
special-token append, a chat-template flag, and load precision --- while being
explicit about which attributions are confirmed and which remain discovery-stage;
and (ii) exhibit the opposite regime, in which the producing instrument of the
original-paper classic models (GPT-J, GPT-NeoX, OPT) is \emph{non-identifiable}
from the surviving public and repository evidence. From this contrast we derive a
minimal, backward-compatible provenance standard --- three recorded fields
(\texttt{model\_revision}, \texttt{torch\_dtype}, \texttt{transformers\_version})
account for every disagreement we attribute, and cover most of the
deployment-side gap at negligible cost --- and an artifact-release checklist that
addresses the engine and hardware side it leaves open.
\end{abstract}

\section{Introduction}
\label{sec:intro}

Large-scale benchmark suites such as \helm \citep{liang2023helm} have become the
shared substrate on which the community compares language models. A central
promise of these suites is transparency: for many (scenario, model) pairs the
public corpus ships the exact prompts sent to the model, the completions it
returned, and the metrics computed from them. This apparent completeness invites
a natural reproducibility question --- \emph{re-run the evaluation locally and
check that the numbers match} --- and a natural expectation that open-weight
models, whose weights anyone can download, should reproduce cleanly.

Our starting point was exactly this program: reproduce a large slice of public
\helm on local open-weight hardware. It did not survive contact with the data,
and \emph{the way it failed is the contribution of this paper.} Two facts
dominate. First, a large fraction of the open-weight corpus was originally served
not by the harness itself but through a hosted inference provider (frequently
Together AI), whose software and hardware stack the public record does not
describe. Second, and more fundamentally, the artifact that \emph{does} survive
--- the resolved \file{run\_spec.json} --- captures \emph{what} to generate
(scenario, prompts, decoding parameters, metrics) and the \emph{name} of the
model and its deployment, but almost nothing about the \emph{execution substrate}:
the precision the weights were loaded at, the exact checkpoint revision, the
attention kernel, the tokenizer and chat-template versions, and the software-stack
versions. A benchmark recipe, in other words, is not a complete experimental
record.

Once this is taken seriously, ``did it reproduce?'' fractures into three distinct
questions that the literature often conflates: (i) \emph{can the evaluation be run
at all?} (gated models and datasets, revoked download endpoints, packaging
incompatibilities, unavailable judges --- these are eligibility and environment
filters, not reproduction outcomes); (ii) \emph{can a faithful local execution be
reconstructed} when the key execution parameters were never written down and the
original may have been served by an external provider; and (iii) when local and
public numbers differ, \emph{to what is the difference attributable} --- a genuine
recipe disagreement, a controllable execution-substrate parameter, an
artifact-processing effect, or an irreducibly lost provenance?

\paragraph{Thesis.} We therefore reframe benchmark reproduction as a layered
\emph{system-identification} problem. The measured score is a function
$Y = F(R, D, I, A, U)$ of the recorded recipe $R$, the deployment $D$, the
execution instrument $I$, the artifact-interpretation history $A$, and residual
stochastic/numerical influences $U$ (\S\ref{sec:model}). The public record
exposes only a projection of these variables. Controlled reconstruction can
\emph{attribute} an observed disagreement to a specific layer, and can sometimes
\emph{close} it by supplying the missing parameter; but when provenance has been
lost or transformed, the original measurement can be \emph{non-identifiable from
the surviving evidence examined}. This is a more defensible and more scientifically
interesting claim than a blanket ``open-weight \helm is (ir)reproducible.''

\paragraph{Contributions.} Concretely, this paper contributes:
\begin{enumerate}[nosep,leftmargin=1.4em]
  \item \textbf{A measurement model and a disagreement taxonomy}
        (\S\ref{sec:model}). We factor a benchmark result into recipe,
        deployment, execution-instrument, artifact-interpretation, and residual
        layers; classify observed disagreements into six causes; and separate
        that ``why'' taxonomy from three orthogonal reproduction \emph{targets}
        --- artifact reconstruction, procedural reproduction, and claim-level
        robustness --- clarifying what ``reproduction'' should even mean.
  \item \textbf{\textsc{EvalAudit}, a reproducibility-auditing system}
        (\S\ref{sec:system}): exact \file{run\_spec.json} replay (never recipe
        reconstruction), an execution-substrate reconstruction and search tool,
        version-pinned ``era'' scoring containers that freeze the measurement
        instrument, and a framework-free layered comparison core.
  \item \textbf{Attribution case studies on modern open-weight models}
        (\S\ref{sec:olmo}). On OLMo-family models we attribute concrete
        disagreements to a tokenizer special-token append, to a chat-template
        \code{add\_generation\_prompt} flag whose behavior depends on the
        \code{transformers} version, and to load precision --- while being
        explicit that the precision result is at present a configuration
        \emph{discovery} awaiting ordinary-path confirmation, and that one dense
        OLMo-2 divergence remains unresolved.
  \item \textbf{A recoverable-vs-non-identifiable contrast}
        (\S\ref{sec:redpajama}--\ref{sec:classic}). RedPajama-3B demonstrates
        end-to-end execution under declared historical proxy instruments, whereas
        the original-paper classic models (GPT-J, GPT-NeoX, OPT) carry a metric
        class path that resolves in \emph{no} released \helm version and
        triangulates to unreleased pre-v0.1.0 code, making the producing
        instrument non-identifiable from surviving evidence.
  \item \textbf{A minimal provenance standard} (\S\ref{sec:provenance}). We show
        that three recorded fields --- \code{model\_revision},
        \code{torch\_dtype}, and \code{transformers\_version} --- transitively
        close most of the identifiability gap, and give a backward-compatible way
        to record them (outside the run-spec identity string, so official$
        \leftrightarrow$local pairing is preserved), plus an artifact-release
        checklist.
\end{enumerate}

\paragraph{Scope and honesty about evidence.} This paper is deliberately explicit
about the maturity of each empirical claim. Several results rest on
configuration-discovery probes rather than completed end-to-end reproductions;
we mark these throughout and describe the confirmation experiment that would
promote them (\S\ref{sec:olmo-confirm}). The system, the conceptual framework,
and the recoverable-vs-non-identifiable contrast stand independently of whether
any single confirmation succeeds.

This work builds directly on \eee \citep{eee2026}, a normalized artifact
representation that \emph{detects} reproducibility discrepancies but, by its
authors' own statement, cannot separate checkpoint, quantization, precision,
tokenizer, and generation-kernel effects. Our positioning is complementary and
one sentence long: \emph{\eee detects reproducibility gaps; this work studies
their attribution, possible closure, and identifiability.}

\section{Related Work}
\label{sec:related}

\paragraph{Holistic benchmarking and its artifacts.} \helm
\citep{liang2023helm} standardized multi-metric, multi-scenario LLM evaluation
and, importantly for us, publishes per-instance request/response artifacts and
resolved run specifications. We treat those artifacts as the object of study
rather than as ground-truth-by-construction, and we distinguish the released
official artifact (the historical record of what \helm reported) from a local
rerun (a test of procedural stability). Related evaluation frameworks
\citep{eval-harness} share the property that a compact recipe stands in for a
much larger latent execution configuration.

\paragraph{Reproducibility of ML evaluation.} A substantial literature documents
that reported model comparisons can be sensitive to implementation and
environment details \citep{pineau2021checklist,dodge2019show}, and recent work
specifically finds that LLM benchmark scores shift with prompt formatting,
few-shot ordering, and decoding \citep{sclar2024quantifying,mizrahi2024state}.
Our contribution is orthogonal to prompt-sensitivity studies: we hold the
\emph{recipe} fixed by replaying the resolved specification verbatim, and isolate
the \emph{execution substrate} --- precision, tokenizer/template version, serving
engine, kernels, hardware --- as the locus of residual disagreement, then ask
which of those parameters are recoverable.

\paragraph{Serving-stack and numerical nondeterminism.} That different serving
stacks can produce different outputs from the same weights is known in folklore
and in engine documentation: greedy decoding through Hugging Face
\code{transformers} \citep{wolf2020transformers} and through vLLM
\citep{kwon2023vllm} differ in kernels, batching, tokenizer handling, and
stop-sequence semantics; low-precision inference and attention-kernel choice
change logits enough to flip near-ties. We make this quantitative at the level of
benchmark metrics, tie specific flips to named unrecorded parameters, and
formalize which are identifiable from surviving evidence.

\paragraph{Detection vs.\ attribution: the \eeeshort line.} The immediate
predecessor of this work, \eee \citep{eee2026}, introduces a normalized
per-instance representation that makes official-vs-local discrepancies
\emph{visible} across heterogeneous harness versions, and reports case studies of
detected gaps on Pythia-6.9B \citep{biderman2023pythia}, Vicuna-7B, and
Falcon-7B. Its appendix explicitly leaves open whether a residual difference
arises from checkpoint, quantization, precision, tokenizer, or generation-kernel
choices. The present paper takes that open question as its subject: we supply the
machinery to attribute and sometimes close such gaps, and we characterize when
the answer is fundamentally underdetermined.

\paragraph{Provenance and artifact preservation.} Our provenance recommendation
(\S\ref{sec:provenance}) is in the spirit of model and dataset documentation
efforts \citep{mitchell2019modelcards,gebru2021datasheets} but targets the
\emph{execution} rather than the model or data: what must be recorded so that a
benchmark score is reconstructible. The closest operational analogue is
containerized, digest-pinned execution; we extend it with version-pinned scoring
instruments for historical corpora and a per-parameter identifiability analysis.

\paragraph{Efficient benchmarking (peripheral).} A separate thread explored
whether a small, well-chosen coreset of questions plus regression can predict
full-benchmark scores \citep{bowyer2026efficient} and whether model-centric
representations help. That thread is not required for the claims here and is
summarized only as future work (\S\ref{sec:limitations}).

\section{A Taxonomy of Benchmark Reproduction}
\label{sec:model}

\subsection{What produces a score}

A measured benchmark result is produced by five things acting together, and naming them separately is what makes the rest of this section possible. The \textbf{recorded recipe} is the scenario, dataset, prompt construction, in-context examples, decoding parameters, and scorer. The \textbf{deployment} is the weights, the model and tokenizer revisions, the chat template, the load precision, and any quantization. The \textbf{execution instrument} is the harness version, serving engine, package and kernel versions, hardware, device topology, and batching or caching behavior. The \textbf{artifact interpretation} is the storage schema, its migrations, canonicalization, and any carried-forward outputs. What is left over is the \textbf{residual}: stochastic and numerical influences that no record fixes. We call these five the coordinates of a run.

The public HELM record exposes the recipe, and the names but only a projection of the values within the deployment; of the instrument, the artifact interpretation, and the residual it exposes almost nothing. Reproduction is therefore an inverse problem. One chooses a deployment and an instrument, replays the recorded recipe, and asks whether the result matches the published one --- and when it does not, which coordinate is responsible. We say a run is identifiable at a stated tolerance when every materially compatible choice of coordinates produces equivalent results under a declared output-, metric-, or conclusion-level criterion.

\paragraph{The taxonomy these coordinates imply.} The decomposition dictates how reproduction attempts can be classified, and that classification organizes the rest of the paper. Whether a run can be executed at all is a precondition rather than an outcome, and forms an upstream gate (\S\ref{sec:gate}). Past that gate, an attempt is characterized along two orthogonal axes: which coordinates it puts under test (\S\ref{sec:targets}), and, when its result disagrees with the record, which coordinate accounts for the difference (\S\ref{sec:taxonomy}). Both axes follow from the decomposition rather than being proposed alongside it --- the first partitions the coordinates into those an attempt exercises and those it holds fixed, and the second assigns a disagreement to one coordinate, to the residual, or to no coordinate at all. Neither axis refines the other, and collapsing them is what produces the familiar but uninformative verdict that a benchmark ``did not reproduce.''

\subsection{The eligibility gate}
\label{sec:gate}

Before any disagreement is measured, a run must be executable. Gated models or datasets, revoked download access, missing credentials, unsupported packaging, unavailable judges, and infeasible resource requirements are eligibility and environment filters. They answer ``can we run this?'', not ``did it reproduce?'', and must never be folded into a negative reproduction result. Making this gate explicit --- every excluded run carries a typed exclusion reason --- is what allows a defensible denominator: how many runs were runnable, how many were excluded, and for which reason each.

The gate is relative to a target, not to a run. The sharpest case is a model-graded scenario whose judge is no longer served: the judging step is unrepeatable, so the run is ineligible for procedural reproduction, but the judge's per-instance verdicts are themselves part of the released record, and the published aggregate can still be recomputed from them and checked. Such a run is out of scope for the second target of \S\ref{sec:targets} and fully in scope for the first. Treating eligibility as a property of the run alone would discard precisely the evidence that survived, and would understate the auditable fraction of the corpus; we therefore scope every exclusion reason to the target it blocks.

\subsection{Axis 1: what is under test}
\label{sec:targets}

A reproduction attempt must first declare what it is attempting to reproduce. We adapt the methods, results, and inferential trichotomy of \citet{goodman2016does} to a benchmark corpus that publishes per-instance outputs, as described in \S\ref{sec:related}. The three resulting choices each exercise a different subset of the coordinates, and they are not interchangeable:

\begin{enumerate}[nosep,leftmargin=1.6em]
  \item \textbf{Artifact reconstruction}: recompute the published aggregates from the released prompts, completions, and metric records, without rerunning a model. Tests the artifact interpretation and the scorer, and nothing else.
  \item \textbf{Procedural reproduction}: execute the recorded recipe under a pinned or explicitly equivalent instrument, and compare instance-level outputs and metrics. Tests the deployment and the instrument.
  \item \textbf{Claim-level robustness}: determine whether the substantive conclusion --- a ranking, or a model-generation improvement --- survives plausible execution substrates even when byte-identical outputs do not.
\end{enumerate}

Targets (2) and (3) correspond to Goodman's methods and inferential reproducibility, the former with its ``same data and tools'' premise necessarily relaxed, since the original tools are what did not survive; target (1) has no counterpart there, and his results reproducibility has none here, because a fixed historical corpus admits no new study.

Released official artifacts remain the historical ground truth for what HELM reported. A rerun assesses the stability and identifiability of the procedure that produced them and does not overwrite the record, which is what licenses us to report a differing rerun as evidence about the instrument rather than as a correction of the official number.

\subsection{Axis 2: why two executions disagree}
\label{sec:taxonomy}

Once the target is fixed and the attempt executes, an observed disagreement admits one of six causes --- one per coordinate, plus a residual and an underdetermined case:

\begin{enumerate}[nosep,leftmargin=1.6em]
  \item \textbf{Recipe mismatch} --- the recorded recipe differs: prompt construction, in-context examples, decoding, or scorer.
  \item \textbf{Deployment mismatch} --- the deployment's client, endpoint, or engine identity differs, as when a hosted API is replaced by a local one, or in-process \texttt{transformers} by vLLM.
  \item \textbf{Execution-instrument mismatch} --- the instrument differs: precision, tokenizer or template version, attention kernel, device placement, or software-stack versions.
  \item \textbf{Artifact-interpretation mismatch} --- the interpretation differs: row reordering, schema migration, or carried-forward outputs.
  \item \textbf{Residual disagreement under controlled equivalence} --- every controllable coordinate is matched and a gap remains. This is the only category analogous to conventional repeatability failure.
  \item \textbf{Historically non-identifiable reproduction} --- the surviving evidence does not uniquely specify what execution should count as faithful. Qualitatively different from (5): not ``it failed to repeat'' but ``what would count as a faithful execution is underdetermined.''
\end{enumerate}

The two axes are independent: any target can fail for any cause, and it is their combination, not either alone, that describes an outcome. We deliberately avoid the phrase ``true reproducibility failure'' whenever execution provenance is incomplete. Categories (1) through (4) are usually recoverable by controlling the substrate, and (6) is a first-class result rather than a gap to be closed.

\subsection{Per-parameter identifiability}
\label{sec:identifiability}

Identifiability in the sense fixed above --- equivalent results, within a stated tolerance, across every materially compatible configuration --- is not a binary property of an entire run. Each latent parameter within the deployment and the instrument has its own status, and the per-case analysis assigns one to each (Table~\ref{tab:identifiability}). This refinement is what makes both the positive attributions and the non-identifiability result precise: we never claim to identify a whole historical deployment, only specific coordinates, each with its own evidentiary basis.

A coordinate can fail to be recoverable in two ways, and the difference decides what can be done about it. It may be structurally unrecoverable, in that no record of this form could pin it down, or merely practically unrecoverable, in that the particular record which happens to survive is insufficient \citep{raue2009identifiability}. Our classic-model result (\S\ref{sec:classic}) is the second kind, and that is what makes the provenance recommendation of \S\ref{sec:provenance} actionable: a practically non-identifiable parameter is one a richer record would have recovered, whereas a structurally non-identifiable one would leave a recording standard nothing to fix. We therefore state that result as non-identifiability relative to named evidence rather than as a claim of irreproducibility.

\begin{table}[t]
\caption{Per-parameter identifiability in the studied cases. The status is per
coordinate, not per run; \S\ref{sec:cases} assigns each case's disagreement to a
coordinate and to its status here.}
\label{tab:identifiability}
\begin{center}
\small
\begin{tabularx}{\linewidth}{@{}%
  >{\raggedright\arraybackslash\hsize=0.80\hsize}X
  >{\raggedright\arraybackslash\hsize=0.55\hsize}X
  >{\raggedright\arraybackslash\hsize=1.65\hsize}X@{}}
\toprule
\textbf{Parameter} & \textbf{Typical status} & \textbf{Basis / limitation} \\
\midrule
Run specification & identifiable & frozen public \texttt{run\_spec.json};
migration can still alter interpretation. \\
Client / deployment class & often identifiable & recorded deployment name and
harness config; hosted-provider internals remain opaque. \\
Load dtype & conditional & explicit when pinned; else inferable only from loader
source and package-version bounds. \\
Tokenizer special-token behavior & often recoverable & repo/config inspection and
token-level prompt comparison. \\
Chat template / generation prompt & often recoverable & revisioned tokenizer files
or controlled rendering sweeps; effective behavior depends on package version. \\
Model / tokenizer revision & frequently missing & rarely pinned; mutable aliases
prevent unique historical recovery. \\
Engine / package / kernel versions & often missing & client family may be known;
exact versions and kernels usually not. \\
Batching, cache, topology, hardware & usually missing historically & measure
robustness rather than promise exact replay. \\
Harness era & tagged or proxy-only & tagged releases containerizable; migrated
paths may admit only declared proxy instruments. \\
\bottomrule
\end{tabularx}
\end{center}
\end{table}

\section{Methodology}
\label{sec:system}

Every result below is produced by one instrument, \textsc{EvalAudit}, which
operationalizes the taxonomy of \S\ref{sec:model}. It
is an evaluation-auditing and experimental-control system rather than a thin HELM
wrapper: it replays recipes exactly, reconstructs and searches the execution
substrate, freezes the scoring instrument by era, and reports disagreement at the
earliest layer it can localize. We first describe the workflow end to end
(\S\ref{sec:pipeline}), then the four components: three that put a coordinate of
\S\ref{sec:model} under experimental control, and one that assigns a disagreement
to one of them. A per-run \emph{evaluation capsule} (\S\ref{sec:provenance})
records the provenance each component depends on.

\subsection{The workflow}
\label{sec:pipeline}

A reproduction attempt is six stages (Table~\ref{tab:pipeline}), and the split
between them decides what a defect costs to fix.

\begin{table}[t]
\caption{The audit workflow. Commands are prefixed \texttt{eval-audit-}. The
first two stages consume GPU time and produce evidence; the last four are
read-only over artifacts already on disk, so an analysis defect is repaired by
re-rendering rather than re-executing.}
\label{tab:pipeline}
\begin{center}
\footnotesize
\begin{tabularx}{\textwidth}{l l >{\raggedright\arraybackslash}X}
\toprule
\textbf{Stage} & \textbf{Entry point} & \textbf{What it produces} \\
\midrule
Selection     & \texttt{index-historic}           & An inventory of the public corpus with the typed eligibility gates of \S\ref{sec:preconditions} applied and a reason recorded for every excluded run. \\
Execution     & \texttt{run}                      & A local run per eligible official run, replayed from the resolved official spec (\S\ref{sec:fromspec}), writing the same artifacts HELM itself writes. \\
Normalization & \texttt{prepare-eee}              & Both sides rewritten into one framework-neutral per-instance representation \citep{eee2026}, on demand. \\
Composition   & \texttt{build-virtual-experiment} & A declared slice resolved into \emph{packets}: one official row and the local runs claiming to reproduce it, paired at three strictnesses (\S\ref{sec:diff}). \\
Comparison    & \texttt{analyze-experiment}       & Per packet: agreement curves over the tolerance sweep, comparability facts, and a layered diagnosis. \\
Aggregation   & \texttt{build-summary}            & Corpus-level views: the coverage funnels, per-metric drift, and quantile-bucketed examples. \\
\bottomrule
\end{tabularx}
\end{center}
\end{table}

\emph{Selection} walks a mirror of the public corpus, reads each run's recorded
spec, and applies the preconditions of \S\ref{sec:preconditions} as typed gates:
structural completeness first, then modality, weight availability, size, and
local deployability per model. Every excluded run keeps its reason, which is what
makes a typed exclusion a reportable outcome rather than a silent drop.

\emph{Execution} produces the local side. Runs are scheduled one run-entry at a
time across GPU hosts and write the same four files an official run does --- the
resolved spec, the frozen scenario state, the aggregated statistics, and the
per-instance statistics --- so official and local are the same shape on disk
before anything compares them. Model serving is out-of-process behind a local
endpoint, which is what makes the serving half of the execution instrument a
variable we can set (\S\ref{sec:deploymatch}) independently of the harness half
(\S\ref{sec:eras}). The harness half can additionally run inside a container
image pinned by digest and recorded per run, so the software environment is an
audited constant rather than whatever the host happened to have installed.

\emph{Normalization} converts both sides into the normalized per-instance
representation of \citet{eee2026} before any comparison. The comparison core
therefore never reads a HELM
directory; it reads two normalized runs. That is what makes the layered diff of
\S\ref{sec:diff} framework-free rather than HELM-specific, and it is why the same
core can compare an official against a replay, a replay against a repeat, or
either against a run from a different harness entirely. Official runs are
converted on demand, because converting the whole public corpus is the slow step.

\emph{Composition} is where comparisons stop being enumerated by hand. A
\emph{virtual experiment} is a manifest declaring a scope over two indexes ---
one of local audit results, one of official public runs --- which the composer
resolves into packets and into the coverage funnel of \S\ref{sec:coverage}. The
manifest, not the resulting directory, is the archived unit: re-running it
re-derives the slice. \emph{Comparison} then analyses each packet independently
and \emph{aggregation} reads those per-packet reports back to render the
corpus-level views, including quantile-bucketed best/worst/flagged examples so
that extremes are inspected rather than averaged away.

The substrate search of \S\ref{sec:deploymatch} is deliberately not a stage in
this line. It is invoked when a comparison diagnoses a divergence the recorded
evidence does not explain, and what it returns is a configuration change to
execution --- so an attribution closes a loop back through the pipeline rather
than terminating it. Each generated directory records the command that produced
it and stamps its own history, which is what allows a stale report to be
re-rendered from surviving artifacts; \S\ref{sec:limitations} depends on that
distinction when it prices the outstanding work.

\subsection{Exact recipe replay, not recipe reconstruction}
\label{sec:fromspec}

Public HELM runs do not record the CLI string that produced them, but they do
ship a fully resolved \texttt{run\_spec.json}: the complete adapter spec, scenario
spec, metric specs, and annotators, with every default already materialized. The
prior approach --- taken by essentially every earlier reproduction, including our
own first attempts --- \emph{reconstructs} a run-entry string from the run and
re-parses it. Re-parsing re-derives the recipe under \emph{today's} library
defaults, which silently drifts from the recipe that produced the official
numbers.
\par\smallskip\noindent\textbf{Finding.}~\emph{Never reconstruct the recipe; replay the resolved \texttt{run\_spec.json}
verbatim.}\par\smallskip
Our materializer deserializes the official spec and drives an in-process
benchmarking call; the reconstructed run-entry string is used only to \emph{locate}
the official directory, never to drive execution. Two failure modes are handled
asymmetrically because they demand opposite treatment: a scenario/metric
\emph{class} that does not resolve under the local harness is caught loudly at
preflight (an environment mismatch, surfaced as such), whereas a dropped
\emph{field} is silent (the codec fills defaults). We guard the silent mode with a
serialize$\to$deserialize$\to$serialize round-trip test that exploits a
writer/serializer asymmetry: because HELM writes the spec dropping only
\texttt{None} values, any key missing after a round trip is an unambiguous silent
drop.

Because a local replay serves the model on different hardware than the original,
a naive by-name replay would record the \emph{official} deployment identity and
report \texttt{same\_deployment=yes} --- masking the single most important
difference the audit exists to surface. We therefore rewrite the deployment
identity to the local endpoint via an explicit argument (a wrong value fails
loudly \emph{before} any instance runs), so the comparison correctly reports the
substitution.

\subsection{Substrate reconstruction: sweep, don't guess}
\label{sec:deploymatch}

When an unrecorded execution parameter is suspected, the temptation is to guess a
single knob. Two early single-knob theories about the same divergence were both
wrong (\S\ref{sec:olmo}), which motivated a disciplined alternative. Given any
public run, our deployment-search tool builds a grid of plausible local
deployments for that model, evaluates a length-stratified prompt sample, and ranks
each cell by agreement with the official completions (Together stores no
log-probs, so the objective is completion-text agreement plus a model-agnostic
prompt-independence-collapse detector). The grid is two-tier: serve-time knobs
(dtype $\times$ tokenizer $\times$ \texttt{max\_model\_len}) become distinct
endpoints, while request-time knobs (\texttt{add\_special\_tokens},
\texttt{add\_generation\_prompt}, protocol) are probed per request.

A design principle distinguishes two opposite prunings. \emph{Feasibility
pruning} (safe) drops a configuration that physically cannot serve on the hardware
with a typed reason (e.g.\ float32 on a mixture-of-experts model exceeds the vLLM
fused-kernel shared-memory limit). \emph{Relevance pruning} (``this knob probably
won't matter'') is refused: it is precisely the guess the tool exists to avoid.
The tool has \texttt{dry-run}/\texttt{run}/\texttt{confirm} phases; \texttt{confirm}
promotes a winning cell into a real, servable one-endpoint configuration for the
ordinary execution path (the confirmation step of \S\ref{sec:olmo-confirm}).

\subsection{Era-pinned scoring instruments}
\label{sec:eras}

A large fraction of any historical corpus was produced by an \emph{older harness}:
the measurement instrument itself has versions. Concretely, the exact
\texttt{pandas} version flips per-instance row ordering on some scenarios, so the
harness's own dependency versions change the recorded instances. Because the
modern runner pins a recent HELM and imports module paths that do not exist in
older releases, it cannot replay these runs at all.

\textsc{EvalAudit} freezes the harness half of the execution instrument
(\S\ref{sec:model}) by era: each era gets a CPU-only
container whose HELM harness is checked out at that era's release commit, with
era Python and era dependency pins, while model inference stays out-of-process on
a modern serving lease. One manifest maps to one era maps to one image maps to one
measurement instrument, enforced at manifest time and guarded by an
image$\leftrightarrow$era label check. A standalone shim provides a
flag-compatible replay CLI plus a backported serving client, so the era image
depends only on era HELM. This isolates the scoring code as a controlled
variable, separable both from the model and from the serving half of the
instrument, which stays modern and is varied independently --- a prerequisite for
treating a historical disagreement as evidence about the execution instrument
rather than noise.

\subsection{Layered differential comparison}
\label{sec:diff}

All reproducibility verdicts flow through a single framework-free comparison core
that consumes two normalized runs and produces agreement rows, agreement curves,
tolerance quantiles, a facts-grade diagnosis, and a judge-dependent vs.\
deterministic metric-class split. It localizes the \emph{earliest} divergence
layer along
\[
\text{recipe} \rightarrow \text{prompt} \rightarrow \text{tokens} \rightarrow
\text{generation} \rightarrow \text{text} \rightarrow \text{metric},
\]
which is what turns a bare ``the numbers differ'' into an attribution. To keep
comparisons honest, comparability \emph{facts} (e.g.\ same recipe, same
deployment, same judge) are reported separately from the \emph{diagnosis}: a fact
is never re-labeled to hide a substitution.

Pairing an official row with a local one is itself such a decision, and not a
fact read off the record, since the harness that wrote each side decided what
that side contains. We therefore never report a single coverage number, but join
at three nested strictnesses. \emph{Recipe-identical} is agreement on HELM's own
\texttt{run\_spec\_hash}. \emph{Recipe-canonical} is agreement after we
additionally collapse the fields newer HELM populates and older HELM omitted
entirely --- \texttt{chain\_of\_thought\_prefix} and \texttt{\_suffix},
\texttt{global\_suffix}, \texttt{num\_trials}, \texttt{model\_deployment},
\texttt{metric\_specs}, \texttt{groups}, \texttt{annotators}. \emph{Logical} is
agreement on scenario, model, and adaptation, ignoring \texttt{suite\_version}.
The gaps between the three are what the levels are for: identical-to-canonical
is harness-version churn inside the record, and canonical-to-logical is genuine
recipe drift.

Building this core
\emph{additive-first} --- proven byte-equal to the incumbent on fixtures at
$\texttt{atol}=10^{-9}$ before any renderer was switched to it --- is itself part of
the reproducibility discipline: a looser tolerance is treated as a finding to
investigate, not a knob to widen.

\section{Results}
\label{sec:cases}

The results are two corpus-wide measurements, three case studies, and one
extension. The first measurement comes before everything else, because every
comparison depends on it: what it costs to decide that two runs are the same run
at all (\S\ref{sec:coverage}). The second is the broadest reproduction set we
have, 703 paired Qwen comparisons, which establishes the base rate the case
studies are exceptions to (\S\ref{sec:qwen}). The case studies then instantiate
the two regimes of \S\ref{sec:model}. OLMo is
the \emph{partially recoverable} modern case: several disagreements attribute to
named unrecorded parameters (\S\ref{sec:olmo}). RedPajama-3B demonstrates
end-to-end execution under \emph{declared proxy} historical instruments
(\S\ref{sec:redpajama}). The original-paper classic models are the
\emph{non-identifiable} case: their producing instrument cannot be recovered from
surviving evidence (\S\ref{sec:classic}). The extension asks what happens when the
scorer itself is the unavailable component, and substitutes an open judge for a
withdrawn closed one (\S\ref{sec:openjudge}).

\subsection{What pairing costs: harness-version churn in the record}
\label{sec:coverage}

Every official row in scope is joined to local rows at the three nested
strictnesses of \S\ref{sec:diff} --- recipe-identical, recipe-canonical, logical
--- and the gaps between them are the measurement.

\begin{table}[t]
\caption{Pairing an official corpus with its local replays at three nested
strictnesses. The identical column counts agreement on HELM's own run-spec hash;
canonical collapses the schema-evolving fields; logical ignores suite version.
Official rows is the number in scope, not the number attempted.}
\label{tab:coverage}
\begin{center}
\small
\begin{tabular}{lrrrr}
\toprule
\textbf{Experiment} & \textbf{Official rows} & \textbf{Logical} & \textbf{Canonical} & \textbf{Identical} \\
\midrule
Qwen (modern harness)          & 4723 & 775 & 159 & 0 \\
OLMo, six models               &  182 & 149 &  27 & 0 \\
\texttt{gpt-oss-20b} from spec &   11 &   4 &   2 & 0 \\
RedPajama, era \texttt{v0.2.4} &   74 &   2 &   2 & \textbf{2} \\
RedPajama, era \texttt{v0.3.0} &   74 &   2 &   2 & \textbf{2} \\
\bottomrule
\end{tabular}
\end{center}
\end{table}

\par\smallskip\noindent\textbf{Finding.}~\emph{Under a modern harness, byte-identical
recipe agreement is unattainable --- zero in every experiment, including those
whose recipes did not drift at all, because the schema evolved underneath the
hash. It is nonzero in exactly one place: the era-pinned replays, where both
reproduced runs match on HELM's own hash. Byte-identical recipe agreement is not
something a careful modern replay recovers; it is something only an instrument
pinned to the producing era recovers.}\par\smallskip

This measures the instrument where the surviving record makes it visible: the
harness version decides the shape of the file it writes, so two harnesses
describing the same recipe write different bytes. It is a precondition for
everything that follows: a study that paired on
\texttt{run\_spec\_hash} would have found no comparable runs anywhere in this
corpus and concluded, wrongly, that the recipes had drifted. The identical column
is also the sharpest single argument for era-pinned instruments (\S\ref{sec:eras}),
which recover on two runs what no amount of care recovers on 4723.

\paragraph{Scope and freshness.} These counts are read from the derived coverage
layer of the experiment stores, not recomputed for this paper. The raw run
artifacts behind every row survive --- each run path recorded in the stored
comparison manifests resolves on disk --- so the rows are re-renderable from
source rather than needing re-execution. What remains is to re-render them under
current code and hash the inputs against the stored reports
{\color{orange}\textsf{\small[\textsc{pending}: re-render the coverage funnel and
hash its inputs]}}. The end-to-end
\texttt{phi-2} fixtures are excluded: they are tutorial-scope and pair against no
official rows. The Qwen denominator is the full in-scope corpus for that model
family, so its logical column is a coverage figure and not a success rate; most
of the shortfall is eligibility rather than difficulty, since $3282$ of those
$4723$ rows ($69\%$) are vision-language, audio, or omni-modal models that have no
local deployment and are filtered before any attempt.

\subsection{Qwen: the base rate on a modern harness}
\label{sec:qwen}

The Qwen family gives the broadest reproduction set in this study, and therefore
the base rate against which the case studies read as exceptions. Eight
models --- \texttt{qwen1.5-7b}, \texttt{-14b}, \texttt{-32b}, \texttt{-72b},
\texttt{qwen1.5-110b-chat}, \texttt{qwen2-72b-instruct},
and \texttt{qwen2.5-7b/72b-instruct-turbo} --- were replayed from spec and paired
against their official rows, yielding $703$ built comparisons from $775$ planned
run entries with no skipped entries.

Instance-level agreement with the official is measured at zero tolerance, so a
run counts as agreeing only where the per-instance metric values match exactly.
Across the $703$ comparisons the median is $0.990$ and the mean $0.975$;
$608$ ($86\%$) sit at or above $0.95$, $159$ agree exactly, and $10$ fall
below $0.80$, with a minimum of $0.66$. Per-model medians span $0.973$
(\texttt{qwen2.5-7b-instruct-turbo}) to $1.000$ (\texttt{qwen1.5-110b-chat}).

\par\smallskip\noindent\textbf{Finding.}~\emph{On a modern harness with a
from-spec replay, most historical Qwen rows reproduce at or near the instance
level: the median run agrees with the official on $99\%$ of instances at zero
tolerance. Reproduction failure is not the common case; it is the tail, and the
case studies below characterize that tail.}\par\smallskip

This is what makes the rest of the section interpretable. A study that reported
only the disagreements would imply that historical benchmark reproduction is
generally broken, and this corpus says otherwise. What the tail costs to explain,
rather than how large it is, is the subject of \S\ref{sec:olmo} onward.

\paragraph{Scope and freshness.} The counts and agreement values are read from the
store's derived reproducibility rows; all $1406$ run paths recorded in the
comparison manifests resolve on disk, so the set is re-renderable and hashable
without re-execution. These are \emph{reproductions} of rows HELM already
published. The separate arm promised as new Qwen 3.5 evaluations is not reported
here: $72$ runs of \texttt{qwen3.5-9b-base} covering \texttt{mmlu} ($57$),
\texttt{legalbench} and \texttt{wmt\_14} ($5$ each), and one run each of
\texttt{boolq}, \texttt{commonsense}, \texttt{gsm}, \texttt{med\_qa} and
\texttt{narrative\_qa} exist as raw artifacts with per-instance statistics, but
they have never been passed through the analysis and reporting stages, and no
store exists for them. A further nine run directories --- seven \texttt{math} and
two \texttt{natural\_qa} --- were created but are empty, so those two scenarios
were attempted and produced nothing. Because HELM published no Qwen 3.5 rows there is no official
side to pair against, so that arm needs a reporting path of its own rather than
the agreement machinery used here.
\par\smallskip\noindent\fcolorbox{orange}{yellow!12}{%
  \begin{minipage}{0.97\linewidth}\small\textsf{\textbf{[SCAFFOLD --- to be completed]}}\ Run the analysis and reporting stages over the existing
\texttt{qwen3.5-9b-base} artifacts and insert the resulting per-benchmark table as
new evaluations, with the caveat that a single base model is narrower than the
contribution claimed in \S\ref{sec:intro}. No new execution is required.
{\color{orange}\textsf{\small[\textsc{pending}: analyse and report the Qwen 3.5 runs]}}\end{minipage}}\par\smallskip

\subsection{OLMo: attributing modern disagreements to unrecorded parameters}
\label{sec:olmo}

We study six OLMo-family models
\citep{groeneveld2024olmo,olmo2_2024,muennighoff2024olmoe} --- four recent instruct models
(\texttt{olmoe-1b-7b-0125-instruct}, \texttt{olmo-2-1124-7b/13b-instruct},
\texttt{olmo-2-0325-32b-instruct}) on capability benchmarks
(\texttt{bbq}, \texttt{gpqa}, \texttt{ifeval} \citep{zhou2023ifeval},
\texttt{mmlu\_pro} \citep{hendrycks2021mmlu}) and two older base
models with broad classic-track coverage. Three attributions and one open puzzle
emerge.

\paragraph{(a) A tokenizer special-token append (confirmed).} During from-spec
reproduction of \texttt{narrative\_qa}, local OLMo-7B emitted the same
prompt-independent boilerplate for \emph{every} instance while the official
conditioned correctly, with \texttt{prompts\_equal=True} at the string level. The
first hypothesis --- fp16 numerical instability --- was \emph{refuted} by a
deployment sweep. The cause was a tokenizer issue: \texttt{OLMo-7B-hf}'s tokenizer
appends an \texttt{<|endoftext|>} token to the prompt under vLLM's default
\texttt{add\_special\_tokens=True}, which the original deployment did not do and
which HELM never recorded. Serving the append-free tokenizer and forwarding
\texttt{add\_special\_tokens=false} restores prompt conditioning. This is the
paradigm case: an apparent flat irreproducibility was a single unrecorded
request-time flag. The coordinate that differs is the deployment --- the prompt
tokens themselves (Table~\ref{tab:identifiability}, row ``tokenizer
special-token behavior'') --- while the engine default that supplied it sits in
the instrument (\S\ref{sec:model}); fully recoverable either way.

\paragraph{(b) A chat-template version dependency (confirmed).} HELM always calls
\texttt{apply\_chat\_template(\ldots, add\_generation\_prompt=True)}, but that flag
is a no-op unless the chat template \emph{implements} it. The OLMoE template
shipped with HELM's older \texttt{transformers} ignored the flag --- so the model
was fed the prompt \emph{without} the trailing assistant-turn marker --- whereas a
modern \texttt{transformers}/vLLM ships an updated template that honors it and
\emph{appends} the marker, changing the prompt and the output. Rendering with
\texttt{add\_generation\_prompt=false} on the modern stack moves the output back
toward the official, confirmed empirically. The attribution is to the
\emph{tokenizer/chat-template version} coordinate, whose effective behavior
depends on the (unrecorded) \texttt{transformers} version --- a clean example of a
software-version dependency hiding inside a recipe.

\paragraph{(c) Load precision: a discovery, with a deductive leg (partially
corroborated).} HELM's \texttt{HuggingFaceClient} loads a model with only the kwargs
the deployment config supplies; for the OLMo instruct deployments that is
\texttt{device\_map=auto} and \emph{no} \texttt{torch\_dtype}. Under every
\texttt{transformers} 4.x, an unpinned load defaults to \texttt{float32}, ignoring the
checkpoint's bf16 config (reading the checkpoint dtype is a v5 change). Since the
OLMoE architecture landed in \texttt{transformers}${\sim}$4.45, a January-2025 run
sits squarely in the fp32-default regime. This is a \emph{deductive} argument ---
read off the loader source and the version-dependent default, not fit to outputs
--- and is the stronger leg; its residual is the one parameter it depends on, the
exact \texttt{transformers} version, which the run directory bounds but does not
record.

A separate, small \emph{discovery probe} corroborates it: an oracle-scored
deployment-match sweep over $n\approx 12$ \texttt{ifeval} instances finds float32
candidate cells with the highest agreement to the sampled public completions
(Table~\ref{tab:fp32}). The winning \emph{engine} depends on architecture:
OLMoE is MoE, so vLLM's fused-kernel float32 path is infeasible on workstation
cards and only in-process \texttt{transformers} float32 reproduces the sampled
OLMoE outputs, whereas the dense OLMo-2 samples matched most closely under vLLM
float32.

Two ordinary-path \texttt{ifeval} runs at float32 --- on
\texttt{olmo-2-1124-7b-instruct} and \texttt{-13b-instruct}, through the normal
HELM path rather than the probe --- move the strict accuracy about a third of the
way from the default-precision local result to the official
(\S\ref{sec:olmo-confirm}). Precision is thus a real contributor to this gap and
demonstrably not the whole of it. That is short of the confirmation protocol,
which requires held-out instances and a second benchmark, so we continue to report
(c) as a discovery.

\begin{table}[t]
\caption{Deployment-match \emph{discovery} probe ($n\approx 12$ \texttt{ifeval}
instances, scored against sampled public completions; winning cells additionally
use a probe-only \texttt{add\_special\_tokens=false} knob). Float32 is the
highest-agreement tested precision. \textbf{This is configuration discovery, not a
completed reproduction}: matching a small oracle sample does not uniquely identify
the producing deployment, and the confirmation step of \S\ref{sec:olmo-confirm}
is only partially executed. The deductive dtype argument (text) is the stronger leg.}
\label{tab:fp32}
\begin{center}
\small
\begin{tabular}{lll}
\toprule
\textbf{Model} & \textbf{HF best} & \textbf{vLLM best} \\
\midrule
OLMoE-1B-7B (MoE)   & \textbf{fp32 match 0.971} (quasi 1.0) & fp16 partial 0.494 \\
OLMo-2-7B (dense)   & fp32 partial 0.175 (quasi 0.0)       & \textbf{fp32 match 0.915} \\
OLMo-2-13B (dense)  & fp32 partial 0.324 (quasi 0.0)       & \textbf{fp32 match 0.904} \\
OLMo-2-32B (dense)  & fp16 partial 0.234 (quasi 0.0)       & \textbf{fp32-tp2 match 0.961} \\
\bottomrule
\end{tabular}
\end{center}
\end{table}

\paragraph{(d) An unresolved dense-OLMo-2 Hugging Face divergence.} Rigor demanded
following up an inconsistency: the in-process \texttt{transformers} float32 probe
exactly matched the sampled OLMoE outputs but gave near-zero quasi-match for the
dense OLMo-2 models, even though those officials are \emph{recorded as}
\texttt{HuggingFaceClient} deployments. Tokenizing the exact prompt showed
byte-identical tokens across paths, so it is not a prompt difference; yet the
current-\texttt{transformers} float32 first token disagreed with the official while
vLLM float32 matched. A first-token disagreement on byte-identical prompt tokens
localizes the divergence to the forward pass. We report this as an
\emph{unresolved systematic divergence under prompt-token equivalence} --- direct
evidence of execution-stack sensitivity --- and are explicit that its cause is not
yet isolated: candidates include the model/tokenizer revision, the
\texttt{transformers}/\texttt{torch} version, the attention implementation, device
placement, and generation defaults, \emph{and}, not yet ruled out, a defect in our
own dense-model probe path. We reject only the premature dismissal of the
divergence, not the possibility that our tooling contributes to it.

\paragraph{Store status (honest reporting).} The aggregate OLMo drift report on
disk is \emph{stale}: the base-model \texttt{olmo-7b} classic cells carry an
un-deduplicated local-halving artifact (e.g.\ MMLU public/local $0.295/0.144$),
whereas the corrected value is $0.295/0.287$. The code fix is committed but the
store was not regenerated. The raw run artifacts survive, so a correct report
requires only a re-render under the fixed code rather than a re-execution; until
that is done we do not cite base-model OLMo aggregates as results in this draft.
\par\smallskip\noindent\fcolorbox{orange}{yellow!12}{%
  \begin{minipage}{0.97\linewidth}\small\textsf{\textbf{[SCAFFOLD --- to be completed]}}\ Insert the \emph{regenerated} OLMo aggregate score-drift heatmap
(6 models $\times$ ${\sim}12$ benchmarks) after a fresh run with the current
reporting code. Expected qualitative pattern: base-model classic benchmarks
reproduce near-exactly; instruct-model \texttt{ifeval} shows the largest drift,
consistent with the chat-template and precision findings above.\end{minipage}}\par\smallskip

\subsubsection{Generalization probe: GPT-OSS-20B (candidate)}
\label{sec:gptoss}
As a test that the machinery generalizes beyond OLMo, we replayed the four
ungated-judge public \texttt{gpt-oss-20b} rows (\texttt{bbq}, \texttt{gpqa},
\texttt{ifeval}, \texttt{mmlu\_pro}) from spec, after diagnosing and fixing a
null-content chat crash (a null-vs-empty-string serving difference between the
hosted and local engines; normalizing null$\to$``'' is faithful). The regenerated
report shows local and public headline scores agreeing closely: three of the four
scenarios agree exactly on every core metric, and the fourth (\texttt{bbq})
differs by $0.004$, for a worst-case squared error of $1.6{\times}10^{-5}$ ---
nearly three orders of magnitude (${\approx}800\times$) tighter than the mean
OLMo instruct \texttt{ifeval} drift of $1.3{\times}10^{-2}$. We report this as a
\emph{promising candidate}, not a result: the report layer is current-code and the
run artifacts survive, but they have not been hashed against the stored report, so
run-artifact and comparison provenance must be established (re-render, preserve,
hash, and pass a manifest-based acceptance audit) before the numbers are cited.
\par\smallskip\noindent\fcolorbox{orange}{yellow!12}{%
  \begin{minipage}{0.97\linewidth}\small\textsf{\textbf{[SCAFFOLD --- to be completed]}}\ Promote GPT-OSS to a result after re-rendering and preserving the store;
insert the four-cell public-vs-local headline comparison with a preservation
manifest.\end{minipage}}\par\smallskip

\subsection{RedPajama-3B: execution under declared proxy-era instruments}
\label{sec:redpajama}

Not every historical run is unrecoverable. \texttt{redpajama-incite-base-3b-v1}
\citep{together2023redpajama} was run end-to-end through \emph{both} classic-era instruments (v0.2.4 and v0.3.0
scoring containers, \S\ref{sec:eras}) on a small workstation, exercising a
validation ladder: image sanity, instrument fidelity on a \texttt{pandas}-sensitive
scenario (byte-for-byte instance identity, no model), an end-to-end single run, a
full per-era packet, and a per-scenario data-availability audit (data-fetch/auth
failures are classified as environment SKIPs, not instrument failures). This
demonstrates \emph{procedural reproduction under a declared proxy instrument}
(\S\ref{sec:targets}): the era container reproduces \emph{a} faithful scoring
environment for the classic recipe. We are careful to distinguish this from
reconstructing the \emph{original} producing instrument --- the two later eras
carry byte-identical public artifacts for the shared runs (carried forward, not
re-run), so ``two eras agree'' is one public number measured by two instruments,
not two independent confirmations.

Both eras planned and built two packets each, with no skipped run entries, and the
outcome splits by scenario. On \texttt{mmlu} the replay is exact: all eight core
metrics --- exact, quasi-exact, prefix, and quasi-prefix match, on both the test
and valid splits --- agree with the official at $\Delta = 0$, identically under
v0.2.4 and v0.3.0. On \texttt{synthetic\_reasoning\_natural} they do not agree,
and not in the direction drift would predict: the official reports exactly
$0.0000$ on all six set-match metrics while the local run scores $0.147$ to
$0.184$, again identically under both eras. A published zero across every metric
of a scenario is not a value a working scorer ordinarily produces, so we treat
this as an open question about the official record rather than as a reproduction
failure, and count it as neither pending diagnosis. It is a concrete instance of
why the record must be screened before a disagreement is attributed
(\S\ref{sec:preconditions}).
\par\smallskip\noindent\fcolorbox{orange}{yellow!12}{%
  \begin{minipage}{0.97\linewidth}\small\textsf{\textbf{[SCAFFOLD --- to be completed]}}\ Insert the remaining rungs of the validation ladder --- instrument-fidelity
byte match on the \texttt{pandas}-sensitive scenario, and the per-scenario
environment-SKIP counts --- from the preserved \texttt{era-redpajama} stores, whose
raw run artifacts survive and are directly packageable. Resolve the
\texttt{synthetic\_reasoning\_natural} official zero before the table is
published. {\color{orange}\textsf{\small[\textsc{pending}: diagnose the official all-zero scenario]}}\end{minipage}}\par\smallskip

\subsection{GPT-J, GPT-NeoX, OPT: origin-instrument non-identifiability}
\label{sec:classic}

The original-paper classic models
\citep{wang2021gptj,black2022gptneox,zhang2022opt} are the negative counterpart.
Era-replaying them
fails the shim's class preflight: their \texttt{run\_spec.json} names
\texttt{helm.benchmark.basic\_metrics.BasicMetric}, a class path that resolves in
\emph{no} released HELM version. An investigation over a blob-less clone of the
HELM history establishes that the stored path is a once-migrated hybrid --- a bulk
\texttt{benchmark.}$\to$\texttt{helm.benchmark.} prefix rewrite (applied at the
\texttt{src/helm/} rename) of a run-spec whose metric class was still \emph{flat} at
production time. Reversing the naive migration triangulates the producing code to
an ${\sim}$4-week window of \emph{unreleased pre-v0.1.0} HELM (after scenarios
were nested but before metrics were). The class-path \emph{shape} (flat vs.\
subpackage) is thus a sharper origin fingerprint than release dates.

\par\smallskip\noindent\textbf{Finding.}~\emph{The producing instrument for GPT-J / GPT-NeoX / OPT is
\emph{non-identifiable from the surviving public and repository evidence
examined}: no released version resolves the recorded class, and the triangulated
origin is unreleased.}\par\smallskip
This is category (5) of \S\ref{sec:taxonomy}, not a repeatability failure. Because
building the true origin instrument is infeasible, we treat any later-era run as a
\emph{declared proxy} and make the substitution explicit: a self-verifying
class-path canonicalization remaps the migrated class only when its known
relocation resolves to the same leaf (a genuinely wrong era pin still fails
loudly). We scope the claim carefully --- non-identifiable \emph{relative to
surviving evidence}, not an absolute impossibility statement about hypothetical
private provider archives.

\subsection{Substituting an open judge for a withdrawn one}
\label{sec:openjudge}

The preconditions of \S\ref{sec:preconditions} treat an unavailable judge as
blocking procedural reproduction while leaving the released verdicts auditable.
That is the conservative reading. The stronger question is whether a judge that
can still be served substitutes for one that cannot, and if so under what
criterion --- a declared substitution, made explicit exactly as a proxy era is.

We ran two open judges, \texttt{qwen3.5-27b} and \texttt{qwen3.6-35b-a3b}, over
the released responses of two model-graded scenarios, three replicates each, and
compared them against both official judges. The comparison that matters is not
open-versus-official on its own but open-versus-official measured against the
\emph{official-versus-official} baseline: two closed judges already disagree with
each other, and a substitute is admissible when it disagrees no more than they do.

\begin{table}[t]
\caption{Open-judge substitution against the official-versus-official baseline.
The label scenario reports Cohen's $\kappa$ and exact-agreement rate; the
continuous scenario reports signed and absolute mean difference and Pearson
correlation. Each open arm is three replicates over the same released responses.
Two official judges scored each scenario; \textsc{gpt} denotes the stronger of the
two, and the open arms are shown against it. Against the weaker official the open
arms score lower ($\kappa$ $0.852$/$0.854$ on \texttt{xstest}, $r$
$0.779$/$0.824$ on \texttt{wildbench}) but the comparison to the baseline goes the
same way in both cases. On \texttt{wildbench}, $962$ of the $1000$ responses pair
across all arms.}
\label{tab:openjudge}
\begin{center}
\small
\begin{tabular}{llrrr}
\toprule
\textbf{Scenario} & \textbf{Comparison} & \textbf{$\kappa$ / $r$} & \textbf{signed} & \textbf{within} \\
\midrule
\texttt{xstest}    & \texttt{qwen3.5-27b} vs \textsc{gpt} & 0.936 & $+0.009$ & 98.4\% \\
($n{=}450$, label) & \texttt{qwen3.6-35b} vs \textsc{gpt} & 0.928 & $+0.012$ & 98.2\% \\
                   & \emph{official vs official}      & \emph{0.829} & $-0.030$ & \emph{96.0\%} \\
\midrule
\texttt{wildbench}      & \texttt{qwen3.5-27b} vs \textsc{gpt} & 0.826 & $-0.823$ & 63.6\% \\
($n{=}1000$, continuous)& \texttt{qwen3.6-35b} vs \textsc{gpt} & 0.872 & $-0.731$ & 70.5\% \\
                        & \emph{official vs official}      & \emph{0.936} & $-0.127$ & \emph{91.4\%} \\
\bottomrule
\end{tabular}
\end{center}
\end{table}

\par\smallskip\noindent\textbf{Finding.}~\emph{Judge substitutability is a
property of the metric, not of the judge. On the label scenario both open judges
agree with an official judge \emph{better than the two official judges agree with
each other} ($\kappa = 0.936$ and $0.928$ against a $0.829$ baseline), and are
almost perfectly stable across replicates. On the continuous scenario neither
does: both carry a systematic negative bias of roughly $0.8$ points on the
scoring scale, correlate less well with an official judge than the officials do
with one another, and change $39$--$43\%$ of instance scores between replicates
of the same judge on the same responses.}\par\smallskip

The practical consequence is that a withdrawn judge does not uniformly rule out
procedural reproduction. Where the metric is a label, a declared open-judge
substitution is defensible on this evidence and the run returns to scope. Where
it is a continuous score, it is not: the substitution introduces a bias larger
than the disagreement it is meant to stand in for, and the replicate instability
means a single substituted run is not even self-consistent. Reporting the
official-versus-official baseline is what makes that call decidable rather than
rhetorical, and it is the same move as declaring a proxy era instead of quietly
running a later one.

\paragraph{Scope and freshness.} Both analyses are keyed to a response-set hash,
so the open arms score exactly the responses the officials scored, and no model
was re-run to produce them. Two scenarios is a narrow base for a claim about
label versus continuous metrics generally, and the parse-failure rate on the
continuous scenario ($15.9\%$ for the $27$b arm, $2.0\%$ for the $35$b) is high
enough that some of the bias may be parser rather than judge.
\par\smallskip\noindent\fcolorbox{orange}{yellow!12}{%
  \begin{minipage}{0.97\linewidth}\small\textsf{\textbf{[SCAFFOLD --- to be completed]}}\ Extend the open-judge comparison to the remaining model-graded
scenarios and the four smaller open arms already served, and separate parser
failure from judge disagreement on the continuous scenario before generalizing
the label-versus-continuous split.
{\color{orange}\textsf{\small[\textsc{pending}: remaining model-graded scenarios and parser/judge separation]}}\end{minipage}}\par\smallskip

\section{Toward Auditable Evaluation: A Minimal Provenance Standard}
\label{sec:provenance}

The clearest actionable output of this study is a recommendation to the community:
a benchmark harness should record the \emph{execution substrate}, not just the
recipe and the model name. Our investigation of HELM's own deployment metadata
quantifies the gap: of 148 \texttt{HuggingFaceClient} deployments in HELM's
\texttt{model\_deployments.yaml}, only 19 pin \texttt{torch\_dtype} (all
\texttt{bfloat16}) and only 7 pin a model \texttt{revision}. The parameters that most
change a greedy output are precisely the ones least often recorded
(Appendix~\ref{app:unrecorded}).

\paragraph{Three fields close most of the deployment-side gap.} The high-leverage
observation is that many parameters collapse into one:
\begin{quote}
\texttt{model\_revision} $+$ \texttt{torch\_dtype} $+$ \texttt{transformers\_version}
\end{quote}
The revision fixes the tokenizer, chat template, and generation config \emph{as of
then} (they all live inside the model repository); the \texttt{transformers} version
fixes both the precision default and whether the chat template honors
\texttt{add\_generation\_prompt}; the dtype pins the load precision directly. Each of
the OLMo findings (\S\ref{sec:olmo}) reduces to one of these three fields --- (a)
to the revision and (b) to the \texttt{transformers} version, both confirmed, and
(c) to the dtype, which remains a discovery.

We scope that claim deliberately. Of the parameters in
Table~\ref{tab:identifiability} that are not already identifiable, these three
address the deployment coordinate --- dtype, tokenizer behavior, chat
template, revision --- and leave the instrument coordinate untouched: engine
build and kernel versions, batching, cache, device topology, hardware, and
harness era.
Those are precisely the parameters that make a hosted endpoint nondeterministic
even at temperature zero (\S\ref{sec:related}), and no three data fields can pin
them; they require the container digest and engine record that the checklist below
asks for separately. The evidence for the three-field claim is likewise bounded: it
is the mapping above, on one model family, two of whose three legs are confirmed.
What we claim is that these three fields are the highest-leverage additions
available at negligible cost, not that recording them makes a run reproducible.

\paragraph{Record them out of the identity string.} These fields must live in a
machine-readable provenance block (HELM's existing recipe-facts slot suffices ---
no schema change), \emph{not} in the run-spec identity or the
\texttt{model\_deployment} name. A revision SHA in the pairing key would break
official$\leftrightarrow$local matching, whereas a SHA in the provenance block is
invisible to pairing and fully recoverable. Both weight-loading engines already
accept \texttt{revision} as a data-only field; the only code gap is a bundle
boundary that currently drops it.

\paragraph{A reproduction-target-aware release checklist.} Beyond the three
fields, a paper-grade evaluation artifact should preserve (Appendix~\ref{app:checklist}):
source commit and container digest; dataset and scorer versions; model and
tokenizer commit SHAs; chat-template text/hash and tokenized-prompt hash; dtype,
quantization, engine, attention implementation, and generation config; batching,
cache, tensor-parallel, and numerical/hardware settings; raw requests, token ids,
responses, metrics, and artifact-transformation history; \emph{every}
deployment-search candidate, not only the winner; and all analysis/figure scripts.
Crucially, the release should record, for each experiment, its reproduction
\emph{target} (\S\ref{sec:targets}) and three independent freshness facts ---
run-artifact, comparison, and report freshness --- so that a recent report
timestamp is never mistaken for an acceptance criterion (a lesson learned the hard
way when reports outlived the raw runs they summarized).

\section{Limitations and Threats to Validity}
\label{sec:limitations}

We state the boundaries of the evidence plainly; several are active work items
rather than permanent limits.

\paragraph{Discovery vs.\ confirmation.} The precision result
(\S\ref{sec:olmo}(c)) is at present a configuration \emph{discovery} on a small
oracle-scored sample plus a deductive source-level argument, partially corroborated
by two ordinary-path runs (below); it has not been confirmed on \emph{held-out}
instances. Matching a small
sample does not uniquely identify a producing deployment, and the deductive leg is
conditional on the unrecorded \texttt{transformers} version.
{\color{orange}\textsf{\small[\textsc{pending}: held-out OLMo confirmation]}}

\paragraph{The confirmation experiment: begun, not completed to protocol.}
\label{sec:olmo-confirm}
Two ordinary-path \texttt{ifeval} runs at float32 exist for
\texttt{olmo-2-1124-7b-instruct} and \texttt{-13b-instruct}, executed through the
normal HELM path and producing standard artifacts rather than probe output. They
support the attribution's direction without closing it: float32 moves the strict
accuracy from $0.7911$ to $0.7597$ against an official $0.6929$ on the 7B, and
from $0.8555$ to $0.8121$ against $0.7298$ on the 13B --- about a third of the
drift in each case. Precision is therefore a real contributor to the
\texttt{ifeval} gap and not the whole of it, which is a sharper claim than
\S\ref{sec:olmo}(c) makes and a weaker one than a completed confirmation would.

They do not discharge the protocol. Both are \texttt{ifeval} only, on the full
instance set rather than held-out instances, with no second benchmark; neither
has been through the comparison pipeline, so no pair report exists; and the
container ran by local image id rather than a pinned digest, so the runs are not
reproducible across machines. We therefore leave the confirmation open and do not
promote (c) beyond a discovery.
{\color{orange}\textsf{\small[\textsc{pending}: analyse the two float32 runs; extend to held-out
instances, a second benchmark, and a digest-pinned image]}}

The full protocol
separates discovery from held-out confirmation: freeze the discovery instances
and the winning candidate configuration; determine whether the win requires the
probe-only \texttt{add\_special\_tokens=false} behavior, and if so implement it
through the ordinary tokenizer and client path; execute the exact public
\texttt{run\_spec.json} through the normal HELM path rather than the probe,
generating standard artifacts; compare through the standard pair-report
pipeline; evaluate on \emph{held-out} instances and, preferably, a second
benchmark; and preserve the full candidate surface rather than only the winner.
A success would support a concrete attribution-and-closure claim; a failure would
itself be publishable, since it would show that output-matching probes do not
automatically translate into faithful benchmark execution.
\par\smallskip\noindent\fcolorbox{orange}{yellow!12}{%
  \begin{minipage}{0.97\linewidth}\small\textsf{\textbf{[SCAFFOLD --- to be completed]}}\ Insert the confirmation results table: for one dense and one MoE model
$\times$ two benchmarks, report \{naive, exact-recipe, reconstructed-deployment\}
agreement on held-out instances, with exact-token, normalized-text, first-divergent-token,
common-prefix-length, and task-metric columns. {\color{orange}\textsf{\small[\textsc{pending}: held-out confirmation run]}}\end{minipage}}\par\smallskip

\paragraph{Store freshness and preservation.} Some result stores predate the code
that should have produced them: the OLMo aggregate report is stale, and the
GPT-OSS store is a candidate pending a manifest-based acceptance audit. The raw
run artifacts survive in both cases --- every run path recorded in the stored
comparison manifests resolves on disk --- so these are re-render-and-hash tasks
rather than re-executions. What has not been established is that the surviving
artifacts are byte-identical to the ones that produced the stored reports. Until
that audit is done we report these as candidates, not results, and cite no number
that rests on an unaudited store. {\color{orange}\textsf{\small[\textsc{pending}: re-render, hash, and audit the OLMo and
GPT-OSS stores]}}

\paragraph{The audit tool is itself a source of disagreement.} Every number here
passes through our own comparison layer, which decides which stored rows
constitute a run, how completions are normalized before they are compared, and at
what granularity metric keys are matched. Those decisions have been wrong: an
un-deduplicated local join once halved an aggregate that a correct join
reproduces. They are not a coordinate of the original run --- they are a threat to
the validity of our measurement of it --- so we control them by building the
comparison core additively and proving it byte-equal to the incumbent on fixtures
before switching any renderer to it (\S\ref{sec:diff}), and we report a
disagreement only where that core and its predecessor agree.

\paragraph{The exclusion profile is not reported.} \S\ref{sec:preconditions} makes
a typed exclusion a first-class outcome of the taxonomy, but this draft does not
populate it: no corpus-wide breakdown of which runs were excluded, and for which
reason, appears in \S\ref{sec:cases}, and the case studies state only their own
exclusions. The machinery exists and is what produced the scoped corpora used
here; what is missing is a pass over a frozen manifest whose output is fit to
publish. A headline ``fraction reproducible'' statistic would require that
manifest with an explicit numerator and denominator, and fresh stores across
families, so we make no population-level prevalence claim.
{\color{orange}\textsf{\small[\textsc{pending}: corpus-wide exclusion-reason breakdown from a pinned manifest]}}

\paragraph{Single-machine measurements.} Reported agreement is
official-vs-local and same-machine; cross-machine variance (which would separate
platform-specific from recipe-specific drift) is scoped but not completed.
{\color{orange}\textsf{\small[\textsc{pending}: cross-machine repeatability baseline]}}

\paragraph{Oracle objective without log-probs.} Because the hosted originals store
no log-probabilities, the deployment-search objective is completion-text agreement
plus a collapse detector, not a likelihood. This is adequate for ranking
candidates but is not a probabilistic identification.

\paragraph{Non-identifiability is evidence-relative.} The classic-model result
(\S\ref{sec:classic}) is non-identifiability \emph{from the surviving public and
repository evidence we examined}; it does not preclude recovery from a private
provider or author archive we do not have access to.

\paragraph{Score validity is out of scope.} We ask whether a published number can
be reconstructed, not whether it measures what it claims. Open corpora invite
contamination, and a benchmark overstates ability whenever its test data has
reached training, which is rarely checked per benchmark
\citep{sainz2023contamination}. That is an independent axis --- a contaminated
benchmark can be perfectly reproducible and an uncontaminated one irreproducible
--- so we treat the released numbers as the historical record of what HELM
reported and make no claim about their construct validity.

\paragraph{Attribution breadth.} The confirmed attributions are on OLMo-family
\texttt{ifeval}/classic scenarios; ``precision is \emph{the} dominant hidden
variable'' is a hypothesis supported within this pilot, not established across
benchmarks or model families. The efficient-benchmarking thread
(\S\ref{sec:related}) is out of scope here and reserved for future work.

\section{Conclusion}
\label{sec:conclusion}

A public benchmark recipe specifies what to compute, not the experiment that
computed it. Treating open-weight HELM reproduction as a layered
system-identification problem --- rather than a pass/fail rerun --- lets us do
three things the pass/fail framing cannot: attribute concrete disagreements to
named, unrecorded execution parameters (a tokenizer append, a chat-template flag,
a load precision); separate those recoverable cases from a genuinely
non-identifiable one (the classic origin instrument); and derive a minimal,
backward-compatible provenance standard that would have made every confirmed
attribution unnecessary. The strongest durable contribution is the contrast
itself: some lost provenance is experimentally recoverable, and some is
underdetermined by what survives --- a distinction that a uniformly positive
``we reproduced it'' narrative would erase.

We have been explicit about what remains: the precision attribution is a
discovery awaiting ordinary-path confirmation, and several stores must be
regenerated and preserved before their numbers are cited. Those are operational
steps on a clearly specified path (\S\ref{sec:olmo-confirm},
\S\ref{sec:limitations}), and the system, the framework, and the
recoverable-vs-non-identifiable result stand independently of any single
confirmation's outcome. If the community adopts even the three-field provenance
recommendation, the deployment-side disagreements this paper had to
reverse-engineer become, for future evaluations, simply readable off the record.

None of this is specific to HELM. Any suite that identifies an evaluation by a
short specification --- a task name and its configuration, as in
\citet{eval-harness} and \citet{srivastava2023bigbench} --- records the recipe
while leaving the instrument that executed it undescribed, and so inherits the
same gap between what was written down and what was run. HELM is distinguished
not by having the problem but by having published enough per-instance evidence
to make it measurable: a suite that records less does not thereby have a smaller
gap, only less means of seeing it.

\paragraph{Reproducibility of this paper.} \textsc{EvalAudit}, the era containers,
the deployment-search tool, and the analysis pipeline are released, along with the
frozen run specifications, deployment-search candidate surfaces, and figure
scripts. \par\smallskip\noindent\fcolorbox{orange}{yellow!12}{%
  \begin{minipage}{0.97\linewidth}\small\textsf{\textbf{[SCAFFOLD --- to be completed]}}\ Insert the artifact DOI / repository URL and the commit/tag +
container digests used for the camera-ready run. {\color{orange}\textsf{\small[\textsc{pending}: artifact release
bundle, incl.\ a git bundle or immutable reference for the code state]}}\end{minipage}}\par\smallskip


\appendix
\section{The HELM Gotchas Catalogue (condensed)}
\label{app:gotchas}

A reviewer-anticipating ledger of undocumented or hard-to-see HELM behaviors
encountered while building the audit. Each is a concrete pitfall a reproduction
paper's appendix should address; we list the load-bearing ones.

\begin{description}[leftmargin=1em,itemsep=2pt]
  \item[G1. \file{run\_spec.json} schema evolves silently.] Newer HELM populates
    \code{adapter\_spec} fields older HELM omitted, so byte-for-byte hashing yields
    $0/N$ matches across releases even for identical recipes. Workaround: a
    schema-collapsed \emph{canonical recipe hash}.
  \item[G5. \code{instructions} differ across releases on identical scenarios.]
    E.g.\ MMLU prepends ``Answer with only a single letter.''; a genuine recipe
    drift (a different prompt), surfaced via a comparability fact rather than
    silently normalized.
  \item[G7. Same weights $\neq$ same output across serving stacks.] Even greedy,
    \code{transformers.generate()} and vLLM differ (kernels, batching,
    stop-sequence handling, EOS/pad) --- a few-percent per-instance flip rate on
    multiple-choice, tens of percent on long-form.
  \item[G10. Public-track \code{suite\_version} is not the HELM release version.]
    The same spec can appear identically in \code{v0.2.4} and \code{v0.3.0}; era
    replay keys the instrument on \code{(public\_track, suite\_version)}.
  \item[G13. Classic \code{class\_name} paths resolve in \emph{no} HELM version.]
    \code{helm.benchmark.basic\_metrics.BasicMetric} is a once-migrated hybrid
    (a naive \code{benchmark.}$\to$\code{helm.benchmark.} rewrite of a still-flat
    production path), pinning the producing code to unreleased pre-v0.1.0 HELM.
    Class-path \emph{shape} is the origin fingerprint (\S\ref{sec:classic}).
\end{description}
\scaffold{Expand to the full G1--G13 table for the camera-ready, with a
one-line repro pointer for each.}

\section{The Unrecorded-Parameter Taxonomy}
\label{app:unrecorded}

HELM serializes the recipe and the model name but almost nothing about \emph{how}
the numbers were computed. Parameters ranked by effect on a greedy (temperature-0)
output:

\begin{table}[h]
\centering\small
\begin{tabularx}{\linewidth}{@{}p{0.05\linewidth}p{0.30\linewidth}Y@{}}
\toprule
\textbf{Tier} & \textbf{Parameter} & \textbf{Why it matters} \\
\midrule
1 & Model weight revision (SHA) & Same name $\neq$ same weights; fixes tokenizer,
    template, generation config as-of-then. Highest leverage. \\
1 & Load precision / dtype & Unpinned $\Rightarrow$ \code{transformers}-version
    default (pre-v5 $=$ float32). \\
1 & Quantization & GPTQ/AWQ/fp8/bnb change weights outright. \\
1 & Software-stack versions & Sets the fp32-vs-bf16 default \emph{and} chat-template
    behavior; nowhere in the run dir. \\
2 & Attention implementation & eager/sdpa/flash/paged change reduction order
    $\Rightarrow$ different logits at temp 0. \\
2 & Hardware / device topology & determines kernels, TF32, and whether fp32 fits;
    multi-GPU reduction order differs. \\
2 & Batch composition & kernels are not batch-invariant. \\
3 & Tokenizer $+$ chat-template version & the \code{add\_generation\_prompt} drift. \\
3 & \code{add\_special\_tokens} / BOS & the OLMo-7B append case. \\
3 & \code{generation\_config.json} defaults & injects unspecified fields. \\
\bottomrule
\end{tabularx}
\caption{Unrecorded execution parameters (of 148 \code{HuggingFaceClient}
deployments, 19 pin dtype, 7 pin revision). Several Tier-3 gaps collapse into
Tier-1: pinning the model revision recovers the tokenizer, template, and
generation config together.}
\label{tab:unrecorded}
\end{table}

\section{Publication-Grade Preservation Checklist}
\label{app:checklist}

A paper artifact should preserve: source commit and dirty state; container
digests; Python/\code{transformers}/\code{torch}/CUDA/driver versions; GPU model
and topology; model and tokenizer commit SHAs; effective dtype, quantization,
attention backend, TF32/matmul flags, batching, caching, tensor parallelism; raw,
rendered, and tokenized prompts; complete generation configuration; frozen public
run specs and hashes; raw local and public outputs; \emph{every} deployment-search
candidate (not only the winner); compatibility rewrites; report-generation
version; seeds and repeated-run manifests; and all analysis/figure scripts. Each
experiment additionally records its reproduction \emph{target} and its
run-artifact / comparison / report freshness (\S\ref{sec:provenance}).

\section{Claim--Evidence--Status Ledger (excerpt)}
\label{app:ledger}

The camera-ready should carry a single machine-readable claim ledger as the
source of truth for every empirical statement; prose tables are snapshots of it.
An excerpt of the current statuses:

\begin{table}[h]
\centering\small
\begin{tabularx}{\linewidth}{@{}p{0.42\linewidth}p{0.20\linewidth}Y@{}}
\toprule
\textbf{Claim} & \textbf{Status} & \textbf{Required action} \\
\midrule
Reproduction is a layered system-identification problem & established & --- \\
OLMo-7B gibberish caused by tokenizer EOS-append & established & --- \\
\code{ifeval} sensitive to \code{add\_generation\_prompt} & established & --- \\
Historical OLMo dtype most likely float32 & conditional & preserve deployment
    metadata; state version bounds \\
Float32 best-matches sampled OLMo outputs & preliminary (probe) & ordinary-path
    held-out confirmation \\
Dense OLMo-2 HF divergence has a single known cause & unresolved & test substrate
    axes \emph{and} the probe-artifact hypothesis \\
GPT-J/NeoX/OPT identify a unique released producing harness & false & call later
    eras declared proxies \\
Classic origin interval is unreleased pre-v0.1.0 & established (scoped) & preserve
    source-history search scope \\
OLMo/GPT-OSS aggregates are publication-ready & not yet & regenerate + preserve +
    acceptance-audit \\
Most residual gaps are small and attributable & hypothesis & frozen denominator,
    fresh stores, confirmation, ablations \\
\bottomrule
\end{tabularx}
\caption{Excerpt of the claim--evidence--status ledger. Empirical claims are
scoped to what has actually been run; several are explicitly pending the
experiments of \S\ref{sec:limitations}.}
\label{tab:ledger}
\end{table}


\bibliographystyle{plainnat}
\IfFileExists{references.bib}{\bibliography{references}}{%
  {\color{orange}\textsf{\small[\textsc{pending}: run bibtex --- references.bib present but bibliography backend not configured in fallback mode]}}}

\end{document}
